\hypertarget{o10prolong:resolvent-torsion-closes-the-first-actual-prolongation-test}{%
\section{Resolvent torsion closes the first actual prolongation
test}\label{o10prolong:resolvent-torsion-closes-the-first-actual-prolongation-test}}

2026-09-29. This is a further consequence at an arbitrary \textbf{fixed
finite} Euler-density bound. It strengthens the derivative information
about the regularized torsion. It does not, by itself, close general
O10.

\hypertarget{o10prolong:inputs-and-exact-scope}{%
\subsection{1. Inputs and exact scope}\label{o10prolong:inputs-and-exact-scope}}

Consider an actual endpoint sequence with scalar curvature −4, beta
tending to infinity, endpoint defect tending to zero, and a fixed bound
on chi/A. Use the manuscript's actual positive monopole Phi, its cubic
torsion tau in E = S− tensor L, rho = \textbar Phi\textbar², and the
resolvent section v constructed in \texttt{O10\_Critical\_Density.md}.

All norms below use normalized volume. The finite-density bounds and
resolvent construction give

\[
 \|v\|_\infty+\|\nabla v\|_4+\|\nabla^2v\|_2\le C_B,
 \qquad v-\tau\longrightarrow0\text{ in }H^1,
 \qquad D^-_{3A}v\longrightarrow0\text{ in }L^2.
\]

The actual positive-field gradient theorem gives \(|\nabla\Phi|^4\)
uniformly integrable, \(\nabla^2\Phi\) bounded in normalized \(L^2\),
and \(m=|\tau|^2\to m_*=128/81\) in normalized \(L^2\). Consequently
\(v-\tau\to0\) also in normalized \(L^4\): its fourth powers are
uniformly integrable and it converges to zero in measure.

Unlike the small-opposite-Weyl theorem, the present argument also uses
the full strong-current upgrade

\[
 N_\tau=(\nabla\tau)^{1,0}_{J_+}\longrightarrow0\text{ in }L^2.
\]

It follows that

\[
 \mathcal A_v=(\nabla v)^{1,0}_{J_+}\longrightarrow0\text{ in }L^2.
\tag{1}
\]

The almost-complex structure \(J_+\) here is the actual one selected by
Phi. It is used only where rho is bounded away from zero. Low-rho sets
will be removed at the end, not ignored.

\hypertarget{o10prolong:replacing-the-bracket-coefficient-without-changing-the-actual-bracket}{%
\subsection{2. Replacing the bracket coefficient, without changing the
actual
bracket}\label{o10prolong:replacing-the-bracket-coefficient-without-changing-the-actual-bracket}}

The exact actual tensor identity in the manuscript is

\[
 \partial_\nabla\mathcal A_v+
 \mu_{\rm act}(\nabla v)^{0,1}
 =(F^E)^{2,0}v.
\tag{2}
\]

It holds for every section v of E. The coefficient mu\_act is the actual
Nijenhuis tensor of J+. In a compatible positive-spinor frame its
components depend linearly on the conjugate of tau/rho\^{}(3/2); this is
the manuscript's exact cubic bracket identity, not a new almost-complex
structure.

Define the \textbf{comparison tensor} mu\_v by substituting v for tau in
that algebraic formula. This does not assert that mu\_v itself is a
Nijenhuis tensor. Fix a smooth cutoff f(rho), zero for rho at most c and
one for rho at least 2c, with 0 \textless{} 2c \textless{} 4/3. On its
support,

\[
 |\mu_v|\le C_c,
 \quad |\nabla\mu_v|\le C_c(|\nabla v|+|\nabla\Phi|),
 \quad |\nabla J_+|+|df|\le C_c|\nabla\Phi|.
\tag{3}
\]

Here v is bounded, and all coefficients involving rho or the unit
positive spinor are smooth with bounded derivatives on rho at least
c.~Moreover

\[
 \|\mu_{\rm act}-\mu_v\|_{4,\{\rho\ge c\}}\to0.
\]

Therefore

\[
 R=(\mu_{\rm act}-\mu_v)(\nabla v)^{0,1}
 \longrightarrow0\text{ in }L^2\text{ on }\{\rho\ge c\}.
\tag{4}
\]

Equation (2) becomes

\[
 \partial_\nabla\mathcal A_v+ B_v=(F^E)^{2,0}v-R,
 \qquad B_v=\mu_v(\nabla v)^{0,1}.
\tag{5}
\]

The actual curvature identity gives

\[
 |(F^E)^{2,0}v|\le C|\operatorname{Ric}^{\circ}|\,|v|,
\]

so this forcing tends to zero in normalized L2. No opposite-Weyl
derivative or opposite-Weyl smallness enters it.

\hypertarget{o10prolong:the-formerly-missing-adjoint-estimate-is-now-available}{%
\subsection{3. The formerly missing adjoint estimate is now
available}\label{o10prolong:the-formerly-missing-adjoint-estimate-is-now-available}}

From (3), the resolvent budgets and Hölder,

\[
 \|\partial_\nabla^*(f^2B_v)\|_2\le C_{B,c}.
\tag{6}
\]

Indeed differentiation produces only the following types of factors:

\[
 f^2\mu_v\nabla^2v,\quad
 f^2(\nabla\mu_v)\nabla v,\quad
 f\,df\,\mu_v\nabla v,\quad
 f^2\nabla J_+\,\mu_v\nabla v.
\]

The first is L2 bounded; each other factor is a product of two
L4-bounded tensors and bounded coefficients. All
moving-frame/type-projection terms are included in the last type. This
is precisely the estimate unavailable for the unregularized tau in the
old prolongation attempt.

Pair (5) with f² B\_v and integrate over the closed manifold. Since f
vanishes before rho reaches zero, this is a globally legitimate
integration by parts. We obtain

\[
 \int f^2|B_v|^2
 =\operatorname{Re}\int f^2\langle(F^E)^{2,0}v-R,B_v\rangle
 -\operatorname{Re}\int
 \langle\mathcal A_v,\partial_\nabla^*(f^2B_v)\rangle.
\tag{7}
\]

The first term is o(V), because its first factor tends to zero in
normalized L2 and B\_v is L2 bounded. The second is o(V) by (1) and (6).
Thus

\[
 \boxed{\fint f^2|\mu_v(\nabla v)^{0,1}|^2\longrightarrow0.}
\tag{8}
\]

\hypertarget{o10prolong:exactly-which-derivative-has-now-been-removed}{%
\subsection{4. Exactly which derivative has now been
removed}\label{o10prolong:exactly-which-derivative-has-now-been-removed}}

On rho \textgreater{} 0 and v nonzero, choose compatible frames in which
v = sqrt(a) e1, where a = \textbar v\textbar². Put Z1 = e1 − i e2 and Z2
= e3 − i e4, as in the original prolongation section. The algebraic
comparison tensor satisfies

\[
 (\mu_v(\nabla v)^{0,1})(Z_1,Z_2)
 =\frac{\sqrt a}{\rho^{3/2}}\nabla_{\bar Z_1}v.
\]

With the manuscript's tensor norms this reads

\[
 |\mu_v(\nabla v)^{0,1}|^2
 =\frac{a}{4\rho^3}|\nabla_{\bar Z_1}v|^2.
\tag{9}
\]

Both rho and a converge in measure to strictly positive constants. On
the exceptional low-amplitude sets, the L4 bound on grad v gives

\[
 \fint_E|\nabla v|^2
 \le\left(\fint|\nabla v|^4\right)^{1/2}
     \left(\frac{\operatorname{Vol}(E)}V\right)^{1/2}\to0.
\]

Hence (8)--(9), followed by this exceptional-set estimate, imply

\[
 A_1=\nabla_{Z_1}v\to0,
 \qquad A_2=\nabla_{Z_2}v\to0,
 \qquad B=\nabla_{\bar Z_1}v\to0
 \quad\text{in normalized }L^2.
\tag{10}
\]

These assertions mean the invariant corresponding projections on the
region of definition, extended with cutoffs; no globally chosen frame is
required.

Let C = grad\_(bar Z2) v. The original Clifford matrices give, up to an
irrelevant simultaneous Dirac sign,

\[
 D^-v=((A_1)_1-(A_2)_2,\ B_2+C_1).
\]

Thus also C1 tends to zero. The only derivative component not eliminated
is the complex horizontal component C2.

\hypertarget{o10prolong:two-tempting-continuations-do-not-close-the-argument}{%
\subsection{5. Two tempting continuations do not close the
argument}\label{o10prolong:two-tempting-continuations-do-not-close-the-argument}}

\hypertarget{o10prolong:reversing-the-almost-complex-structure-selects-the-wrong-direction}{%
\subsubsection{5.1 Reversing the almost-complex structure selects the
wrong
direction}\label{o10prolong:reversing-the-almost-complex-structure-selects-the-wrong-direction}}

In the ideal remaining configuration, the negative spinor connection has
off-diagonal form

\[
 b=(0,0,c,-ic),
\]

while the actual positive connection has off-diagonal form

\[
 \eta=(\kappa,-i\kappa,0,0),\qquad \kappa\ne0.
\]

For the opposite structure with (1,0) frame (Z1, bar Z2), the Clifford
matrices give

\[
 (\nabla Z_1)^{0,1}_-= -b\,Z_2,
 \qquad(\nabla\bar Z_2)^{0,1}_-= b\,\bar Z_1.
\]

Therefore

\[
 [Z_1,\bar Z_2]^{0,1}_-
 =b(Z_1)\bar Z_1+b(\bar Z_2)Z_2
 =2c\,Z_2.
\]

The positive spinor already has small derivative in the Z2 direction. A
second copy of the same bracket test consequently does \textbf{not}
force c to vanish. Replacing this direction by bar Z1 would be an
incorrect chirality computation.

\hypertarget{o10prolong:the-next-cubic-spinor-lies-in-the-negative-potential-line}{%
\subsubsection{5.2 The next cubic spinor lies in the negative-potential
line}\label{o10prolong:the-next-cubic-spinor-lies-in-the-negative-potential-line}}

The invariant cubic definition is t(Phi) = Clifford(det(Phi,grad Phi))
Phi. For v = sqrt(a) e1 of the opposite chirality, with off-diagonal
derivative components k\_j, the adjoint Clifford matrices give

\[
 t(v)=a(k_1-ik_2,\ k_3+ik_4).
\]

In the remaining configuration k = (0,0,c sqrt(a),−i c sqrt(a)), this is

\[
 t(v)=(0,2a^{3/2}c).
\]

It points along the orthogonal positive-spinor line J Phi, not along
Phi. Its determinant-nine-A Lichnerowicz potential therefore tends to −7
on this line, not +5. The proposed positive-potential vanishing argument
does not apply. These algebraic checks identify the remaining component;
they do not construct an endpoint counterexample.

\hypertarget{o10prolong:research-status}{%
\subsection{6. Research status}\label{o10prolong:research-status}}

The first prolongation estimate (8)--(10) is a genuine new consequence
made possible by resolvent regularization. It works for each fixed
finite Euler density and removes one full additional anti-holomorphic
derivative. General O10 still requires control of C2, or another global
incompatibility using this stronger derivative structure. No such final
incompatibility is claimed in this note.
